\section{Examples and finite verification}
\label{sec:verification}

The proof is combinatorial and applies to arbitrary input size. The programs
included with the supplementary material enumerate finite rotation systems, verify their
face permutations, and test particular identities used in that proof.
They do not implement the complete-sector algorithm or certify the theorem
by experimentation. This section distinguishes these two kinds of evidence.

\subsection{Checks of the new sector and order identities}

\begin{center}
\begin{tabular}{>{\raggedright\arraybackslash}p{0.49\linewidth}r >{\raggedright\arraybackslash}p{0.25\linewidth}}
\toprule
Check & Cases & Scope \\
\midrule
Injective tree-contour transport & 35,240 & Four tree topologies; five or six leaves \\
Self-pipe sector criterion & 720 & Anchor-preserving permutations on three through seven darts \\
Whole common-P order enforcement & 11,824 & Cyclic orders through degree eight, both frames \\
Interval-bundle partitions & 218 & Common P degrees three through seven \\
Restricted gap predicates & 5,564 & All two-value restrictions in those partitions \\
Source-to-common-face map & 4,272 & Degrees three through eight, both copy signs and frame gauges \\
Extra-attachment pullback & 1,755 & Each tested pullback has at most two values \\
\bottomrule
\end{tabular}
\end{center}

The contour checks compare the original leaf gap to its image under the
local tree-corner injection, including independent choices at degree-two
marked nodes. The self-pipe checks enumerate cyclic orders and compare their
actual invariant-category assignments with the proposed orbit-and-sector
criterion. They include nontrivial common orbit lengths and verify both
reference orientations.

The whole-order check uses the economical chain
$\operatorname{or}(a_0,a_i,a_{i+1})=r$ for $1\le i\le k-2$.
Cutting a cyclic order at $a_0$ turns these into a chain of increasing
comparisons, or decreasing comparisons when $r=1$, so they determine the
whole reference order. These are exactly the chain rows used by the
compiler (Lemma~\ref{gen:chain-lemma}). The source-to-face
check independently verifies the reflected-copy convention and the
two-anchor equation, including the fact that a two-element cyclic order
does not constrain its frame parity.

All these checks passed. The underlying proofs are the contour lemma,
the self-pipe inverse, the source-group lemma, and the two-category lemma;
the reported counts are not substitutes for those arguments.

\subsection{Earlier face and signature tests retained}

The first investigation enumerated 1,557 rotation systems of six small
common graphs and 229,372 rotation systems of their attachment-hub
augmentations. The graphs were a four-cycle, $K_4$, the four-rim-vertex
wheel, a three-path theta, a two-sum of two $K_4-pq$ graphs, and the
triangular bipyramid. The 258 attachment-face states checked agree with
the exact auxiliary-hub lemma. Fixed cycle-basis signatures distinguish
the faces in all these embeddings. The two-sum example includes four
different embeddings with changing face eligibility, so this check does
not silently fix one common embedding.

The local signature construction was also checked on 41,885 face-label
assignments on connected incidence patterns with at most four blocks.
These assignments satisfy the connectivity identities used in the proof;
every resulting signature incidence graph is a tree. The general proof
uses those identities explicitly through $U-v$ connectivity, rather than
asserting independent face choices without a consistency argument.

The listed cases directly cover $K_4$, a theta, and a wheel. The repair
proof additionally covers arbitrary collections of noncycle common blocks
meeting at a cutvertex and disconnected common components containing such
vertices. Those unrestricted configurations are covered by the proof, not
claimed to have been exhaustively enumerated by the small-graph programs.

\subsection{A two-pole obstruction that passes common-rotation consistency}
\label{ver:two-pole}

The following instance tests the precise extra constraint retained by the
new compilation. It is not enough to solve two-pole factors separately,
even after requiring compatible common rotations globally.

Let $K$ have the four common paths $p-a-q,p-b-q,p-c-q,p-d-q$. Take
\[
 V=\{p,q,a,b,c,d,w_1,w_2,x,y,r,s,t\}.
\]
In addition to $E(K)$, put the following edges in the two colours:
\begin{center}
\begin{tabular}{p{0.2\linewidth}p{0.7\linewidth}}
\toprule
Type & Edges \\
\midrule
Common & $qw_1,\;pw_2,\;px,\;qy$ \\
Red private & $w_1a,\;w_2c,\;bc,\;da,\;xq,\;py,\;xr,\;ps,\;qt$ \\
Blue private & $w_1b,\;w_2d,\;xs,\;sq,\;pt,\;ty,\;ry$ \\
\bottomrule
\end{tabular}
\end{center}

Both input graphs are connected and planar and the union is biconnected.
The core consists of $K$ together with $w_1,w_2$. Its union is planar: use the order
$a,b,c,d$, place $w_1$ in face $ab$, $w_2$ in $cd$, and the chords $bc,da$
in their corresponding faces. The external factor on $p,q,x,y,r,s,t$,
after addition of a common edge $pq$, is also a positive SEFE instance.
Its nontrivial common component is the path $x-p-q-y$; its rotations
agree automatically and its only face contains all isolated common vertices.
Both colours are planar unions of three pole-to-pole paths and pendant
edges. Thus the common-embedding criterion gives a simultaneous embedding
of that augmented factor.

For explicit global order consistency, the following are planar spherical
rotations of the full graphs; every list is cyclic. The script verifies
them by their face walks.
\begin{center}
\begin{tabular}{c l l}
\toprule
Vertex & Red & Blue \\
\midrule
$p$ & $a,x,s,b,c,w_2,y,d$ & $a,x,b,c,t,w_2,d$ \\
$q$ & $a,d,y,t,c,b,x,w_1$ & $a,d,y,c,b,w_1,s$ \\
$a$ & $p,d,q,w_1$ & $p,q$ \\
$b$ & $p,q,c$ & $p,w_1,q$ \\
$c$ & $p,b,q,w_2$ & $p,q$ \\
$d$ & $p,q,a$ & $p,w_2,q$ \\
$w_1$ & $q,a$ & $q,b$ \\
$w_2$ & $p,c$ & $p,d$ \\
$x$ & $p,q,r$ & $p,s$ \\
$y$ & $p,q$ & $t,r,q$ \\
$r$ & $x$ & $y$ \\
$s$ & $p$ & $x,q$ \\
$t$ & $q$ & $p,y$ \\
\bottomrule
\end{tabular}
\end{center}
Their common restrictions agree at every vertex. In particular, they give
condition (O) for their prescribed block embeddings.

Nevertheless, the full instance has no SEFE. The mixed paths $a-w_1-b$
and $c-w_2-d$, and the red links $bc,da$, force the common P-order
$a,b,c,d$ up to reversal. They put $w_1$ in face $ab$ and $w_2$ in $cd$.
The set $Z=\{x,y,r,s,t\}$ is connected in $U-V(K)$, so it must occupy
one $K$-face $F$. The red pole-to-pole paths through $x$ and $y$ cannot
use faces $bc$ or $da$, since the red crosscuts there separate $p$ from
$q$. Hence $F\in\{ab,cd\}$.

If $F=ab$, consider the common germ $qy$. The red triangle $qw_1aq$
forces it into the corner between $w_1$ and $b$ at $q$, because the
red $p-y-q$ path cannot enter the cap bounded by that triangle.
The blue triangle $qw_1bq$ forces the same germ into the other corner,
between $a$ and $w_1$, because of the blue $p-t-y-q$ path. The common
rotations cannot agree. If $F=cd$, the red triangle $pw_2cp$ and blue
triangle $pw_2dp$ force the common germ $px$ into opposite corners,
using the red $p-x-q$ and blue $p-x-s-q$ paths. This is again impossible.

These arguments take place within one closed face disk and use only
same-colour noncrossing and the common boundary. They permit crossings
between private edges of different colours. In the displayed compatible
rotations, $r$ is in face $ab$ in red and $cd$ in blue, so equality of
common rotations alone indeed fails.

The general compiler retains one variable $t_{K,Z}$ for the whole face of
this component. Every pinned lobe must equal that same value. It therefore
keeps exactly the constraint that the independent-factor shortcut loses.
The example is a negative test of that shortcut, not a hardness reduction.

\subsection{Other refuted simplifications and the parity check}

An eight-vertex theta example already refutes gluing an arbitrary
common-ended two-pole factor independently; a nine-vertex positive example
shows that its common germs need not form one interval among all germs at
the pole. Another positive example distinguishes which colour blocks
contain the common germs, rather than merely recording the set of poles
with common ends. Their full edge lists and rotation certificates are in
the accompanying verification notes and programs.

The four-dart triple-orientation relation has six possible values, so it
is not the projection of an XOR system. This remains a valid obstruction
to the unrestricted affine-reader proposal. The present proof uses complete
categories and injective maps, and converts predicates to affine equations
only after proving that their domain has at most two values.

A separate exploratory enumeration used a theta with a common pendant edge
and one isolated common vertex. It examined 3,413,220 individual rotation
systems and 13,761 pairs of connected planar colour graphs in its prescribed
finite family. Of the pairs, 4,803 admit a SEFE and 8,958 do not; 5,073 of
the negative pairs nevertheless have compatible common rotations. A basic
independently-even linear system happened to agree with the exact results
on this family. This gives no general sufficiency result: the
counterexamples of Gutwenger, Mutzel and Schaefer~\cite{gms2014} to
Schaefer's Conjecture~6.20~\cite{schaefer} already rule one out (see
also~\cite[p.~6]{fkmp2015} and~\cite[p.~137]{fps2017}). That experiment
plays no role in the algorithm.

\subsection{Scope of the checks}

The recorded scripts test finite rotation and contour identities and the
listed counterexample certificates.  They do not implement the entire
algorithm and are not used in place of any general proof.  The source
package contains the test programs and their recorded outputs.  Separate
independent checks compared the face-list criterion of
Section~\ref{rep:section} and the compiled constraints of
Section~\ref{gen:section} with an exhaustive \SEFE{} oracle on small
instances, and compared the exact feasible variable relations before and
after individual engine operations of Section~\ref{cse:section}; they
found no disagreement.  Like the checks above, they are evidence on small
instances, not a proof.
